\documentclass[12pt]{article}
\usepackage[T1]{fontenc}
\usepackage[utf8]{inputenc}
\usepackage{lmodern}
\usepackage{textcomp}
\usepackage{enumitem}
\usepackage{geometry}
\geometry{margin=1in}
\begin{document}
\begin{center}
Answer Key - Health Sciences - Undergraduate Level Assessment 23 \
\small (Answers are ordered exactly as in the question paper.)
\end{center}

\section*{Multiple Choice}
\begin{enumerate}[label=\arabic*.]
\item \textbf{Answer:} D
\\ \textit{Options:} A) The proportion of infected persons willing to be tested; B) The willingness of persons testing HIV positive to begin treatment; C) The availability of resources for persons testing HIV positive; D) All of the above
\\ \textit{Note:} The "test and treat" approach in health sciences relies on the integration of various diagnostic methods, timely treatment protocols, and the collaboration of healthcare providers, making "all of the above" essential components for its effective implementation.
\item \textbf{Answer:} B
\\ \textit{Options:} A) The same dementia; B) Depression; C) A different dementia; D) Social isolation
\\ \textit{Note:} Caregivers for family members with advanced dementia often experience significant emotional and physical stress, leading to feelings of isolation and helplessness, which can result in depression.
\item \textbf{Answer:} A
\\ \textit{Options:} A) Men with frequent unprotected sex; B) Men with infrequent unprotected sex; C) Women with frequent unprotected sex; D) Women with infrequent unprotected sex
\\ \textit{Note:} Pre-exposure prophylaxis (PrEP) is most effective for men with frequent unprotected sex because they are at a higher risk for HIV transmission, and PrEP significantly reduces the likelihood of infection when taken consistently.
\item \textbf{Answer:} C
\\ \textit{Options:} A) first pharyngeal arch.; B) second pharyngeal arch.; C) third pharyngeal arch.; D) fourth pharyngeal arch.
\\ \textit{Note:} The pharyngeal mucosa is innervated by the glossopharyngeal nerve because it derives from the third pharyngeal arch, which is associated with the development of structures innervated by this nerve, including parts of the pharynx.
\item \textbf{Answer:} A
\\ \textit{Options:} A) the fact that they do not have a refractory period.; B) the response of the inner layers of the vagina.; C) having alternating orgasms in different locations.; D) the G-Spot.
\\ \textit{Note:} The correct answer is based on the understanding that women can experience multiple orgasms because they do not undergo a refractory period like men, allowing them to have successive orgasms without needing time to recover.
\end{enumerate}

\section*{True / False}
\begin{enumerate}[label=\arabic*.]
\setcounter{enumi}{5}
\item \textbf{Answer:} True
\\ \textit{Options:} A) True; B) False
\\ \textit{Note:} The prostate gland is anatomically situated at the base of the bladder and encircles the urethra in men, playing a crucial role in urinary and reproductive functions.
\item \textbf{Answer:} True
\\ \textit{Options:} A) True; B) False
\\ \textit{Note:} The refractory period in men can vary significantly due to factors such as age, individual physiological differences, and overall health, leading to durations that can range from a few minutes to up to 24 hours.
\item \textbf{Answer:} False
\\ \textit{Options:} A) True; B) False
\\ \textit{Note:} Cognitive symptoms are not primarily associated with Parkinson's disease, as they are more commonly linked to other neurodegenerative disorders such as Alzheimer's disease, and while they can occur in Parkinson's, they do not define the disease's progression as significantly as its motor symptoms do.
\item \textbf{Answer:} False
\\ \textit{Options:} A) True; B) False
\\ \textit{Note:} Gonorrhea is not transmitted through the placenta during pregnancy; instead, it can be passed to the baby during childbirth if the mother is infected.
\item \textbf{Answer:} False
\\ \textit{Options:} A) True; B) False
\\ \textit{Note:} A vasectomy involves cutting the vas deferens to prevent sperm from being included in semen, but it does not affect the production of testosterone, which continues to be produced by the testes.
\end{enumerate}

\section*{Long Form Response}
\begin{enumerate}[label=\arabic*.]
\setcounter{enumi}{10}
\item \textbf{Answer:} The significance of social support, particularly through close contact, for bereaved individuals is crucial in facilitating their coping mechanisms during the grieving process. Close relationships provide emotional validation, practical assistance, and a sense of belonging, which can significantly enhance an individual's emotional well-being. When bereaved individuals maintain close contact with friends and family, they are more likely to express their feelings and share their experiences, leading to healthier emotional processing. This social support not only alleviates feelings of isolation but also fosters resilience, ultimately aiding in their recovery journey. Therefore, the presence of strong social networks plays a vital role in helping individuals navigate their grief and begin to heal.
\\ \textit{Note:} The answer correctly identifies that close social support offers emotional validation, practical assistance, and a sense of belonging, all of which are essential for bereaved individuals to effectively cope with their loss and improve their emotional well-being during the recovery process.
\item \textbf{Answer:} Interferons play a crucial role in the immune response by serving as signaling proteins that help coordinate the body's defense against viral infections. When a cell is infected by a virus, it produces and releases interferons, which then bind to neighboring uninfected cells. This binding activates pathways within those neighboring cells that lead to the production of antiviral proteins, enhancing their ability to resist viral replication. Additionally, interferons modulate the immune response by promoting the activation of natural killer cells and macrophages, further contributing to the overall defense against infections. Through these mechanisms, interferons effectively create an antiviral state in surrounding cells, bolstering the innate immune response and limiting the spread of viruses within the body.
\\ \textit{Note:} correct because it accurately describes the mechanism by which interferons act as signaling proteins to activate neighboring uninfected cells, leading to the production of antiviral proteins that strengthen the immune response against viral infections.
\end{enumerate}

\vfill
\noindent\textit{End of Answer Key}
\end{document}